\chapter{总结与展望}

\section{研究工作总结}
多元长程时间序列预测需要同时刻画长期趋势、局部波动与跨变量相关关系。针对现有方法在时频协同建模不足、语言模型语义先验难以稳定引入等问题，本文以统一预测骨架构建为基础，以局部多尺度频域增强和全局频域增强为两条互补技术路线，系统研究了冻结大语言模型语义先验与时频建模机制相结合的可行路径。

首先，本文构建了融合时域建模、基础频域表示与冻结语义先验的统一预测骨架。第三章在时域主干之上引入由 Qwen3-0.6B 生成的变量级语义嵌入，并通过轻量适配器、步长感知门控融合与跨模态注意力对齐实现多源信息协同建模。该统一骨架一方面验证了在冻结语言模型主体条件下引入语义先验的可行性，另一方面也为后续不同频域增强模块提供了共享接口与相对统一的比较基础。

其次，本文提出了融合小波包分解的局部多尺度频域增强方法。第四章通过可学习小波包分解、子带独立映射、共享编码器和注意力池化，将原始序列拆解为具有明确频带语义的局部多尺度表示，从而增强模型对峰谷变化、突变片段和复杂局部波动的刻画能力。实验结果表明，该方法在所选公开基准与代表性数据集上能够稳定提升预测性能，并在局部异常、高频扰动和多尺度波动建模方面展现出较强优势。

再次，本文提出了融合频域前馈网络的全局频域增强方法。第五章基于 FreTS 思路，通过复数频谱建模、频域稀疏化和步长感知门控融合，在统一框架下强化模型对全局周期结构与长期趋势信息的建模能力。与第四章的方法相比，该方案更强调从整体频谱中提取长期稳定模式。主实验、消融实验和参数敏感度实验表明，融合频域前馈网络的全局频域增强方法在所选公开基准与代表性数据集上同样能够稳定提升预测性能，并验证了全局频域增强在统一框架中的可行性。

综合全文研究，可得到以下几点结论。第一，冻结大语言模型语义先验并不会自然转化为时序预测收益，只有在轻量适配、跨模态对齐与任务相关融合机制的共同支持下，语义信息才更有可能稳定发挥作用。第二，显式引入频域信息有助于改善长程预测效果，局部多尺度频域建模与全局频域建模在结构机理和建模重点上存在差异：前者更强调局部异常、复杂波动与多尺度细节表达，后者更强调全局周期结构、长期趋势组织与频谱级信息重构；两类方法在所选公开基准与代表性数据集上均取得了较好的综合性能，并体现出较强的互补性。第三，统一预测骨架不仅有助于降低不同增强模块之间的比较偏差，也为后续扩展新的频域模块、语义先验或融合策略提供了较好的结构基础。总体而言，本文在统一骨架下融合语义先验与频域增强策略，获得了较好的预测性能与结构可解释性，为大语言模型在多元时间序列预测中的应用提供了一条可行研究路径。

\section{研究不足与局限性}
尽管本文围绕统一骨架、局部多尺度频域增强和全局频域增强进行了较为系统的研究，但仍存在若干局限。

首先，在数据层面，本文实验主要基于 ETT、Weather、ILI 和 Exchange 等公开基准数据集。尽管这些数据集覆盖了工业负荷、气象、公共卫生和金融等典型场景，但与真实业务环境相比，其数据规模、缺失模式、采样不均匀性和概念漂移程度仍较为有限，因此本文方法在更复杂真实场景中的泛化能力仍有待进一步验证。

其次，在模型层面，本文采用 Qwen3-0.6B 作为冻结语义编码器，重点验证了在冻结语言模型主体条件下引入语言模型先验的有效性，但尚未系统比较更大规模语言模型、领域专用模型以及不同冻结或参数高效微调策略对时序建模的影响。同时，第三章已经通过无语义描述、描述打乱和随机嵌入替换等控制实验对语义先验的作用进行了初步验证，但围绕更丰富提示模板、领域知识强弱以及外部知识来源差异的系统分析仍可进一步细化。另一方面，第四章的小波包增强方法与第五章的 FreTS 增强方法虽然都建立在统一骨架之上，并在损失函数层面保持了相同的 MSE 主损失 + 稀疏正则训练口径，但二者在频域结构设计、部分结构超参数以及实验呈现方式上仍存在差异，因此现有结论更适合说明两类频域建模机制的相对特点，而非穷尽所有配置下的绝对优劣；二者在统一训练框架下的动态组合、联合优化与输入依赖选择机制也尚未展开。

再次，在实验评估层面，本文各组实验均在统一训练流程下采用随机种子 2020--2030 进行了重复实验，从总体趋势上验证了结果的稳定性，但正文目前尚未系统汇报各方法的均值、方差及显著性检验结果，因此对稳定性与可重复性的呈现仍有进一步完善空间。

最后，在评价与应用层面，本文主要采用 MSE 与 MAE 作为核心评价指标，并据此观察模型在不同数据集和预测长度下的表现。然而，真实业务中往往还需要关注峰值误差、区间预测质量、风险成本或下游决策损失，因此仅依赖 MSE 与 MAE 仍不足以完全刻画模型的业务适配性。此外，本文尚未对推理时延、显存占用、参数规模、训练耗时和部署开销进行系统评估，因此当前结论更适用于离线预测与方法比较场景。

\section{未来研究方向}
针对上述局限，未来可从以下几个方面进一步拓展。

首先，可以研究更加自适应的频域增强与模型选择机制。本文已经表明局部多尺度频域建模与全局频域建模在结构设计和建模方式上存在差异，因此后续可进一步设计输入依赖的选择策略，根据频谱熵、波动强度与趋势稳定性，在小波包增强与频域前馈增强之间进行动态切换或加权融合，并进一步探索面向 MSE、MAE、峰值误差等多目标指标的联合优化方法。与此同时，还应补充在完全同损失函数、同训练预算和同调参范围下的严格控制变量实验，以更清晰地区分结构改进与优化目标变化各自带来的收益。

其次，可以进一步探索更大规模或更具领域知识的大语言模型与时序建模的结合方式。未来可在保持统一骨架设计思想的前提下，研究更高参数规模模型、领域适配模型以及参数高效微调策略在时间序列任务中的适用性，并分析语言模型规模、语义质量与训练代价之间的平衡关系。同时，尽管第三章已经通过无语义描述、描述打乱和随机嵌入替换等控制实验对语义先验的有效性进行了初步验证，但后续仍可结合更丰富的文本描述、事件知识和外部检索信息，构建更完整的时序--语义联合建模框架，并围绕提示模板设计、领域知识强弱、嵌入生成方式以及外部知识来源差异开展更细致的对比实验，以进一步识别语义先验的真实贡献。

再次，可以将当前方法推广到更加复杂的数据形式与下游任务。针对缺失值、多粒度采样、异步多变量及图结构时空数据等更复杂场景，可进一步扩展本文提出的统一骨架与频域增强模块；针对异常检测、状态识别、区间预测和决策支持等任务，也可研究预测模块与业务目标的联合优化方式，以增强模型在真实应用中的实用价值。

最后，在工程应用层面，仍有必要围绕效率、稳定性与部署问题展开研究。例如，可通过模型蒸馏、量化、结构剪枝和算子融合等手段降低统一框架在训练与推理阶段的计算开销；也可结合更系统的多随机种子统计展示、显著性检验、基线独立调参与资源消耗评估，更全面地分析模型在资源受限环境下的稳定性、可复现性与可部署性。若能够在保持精度的同时进一步降低推理时延和显存占用，本文提出的时频增强框架将在实际在线预测场景中展现出更强的应用潜力。

总体而言，本文验证了语义先验与频域增强相结合在多元长程时间序列预测中的有效性，也揭示了局部多尺度频域建模与全局频域建模在作用侧重点上的差异。未来若能沿着自适应模型选择、更强语义先验、复杂场景泛化与高效部署等方向持续推进，将有望进一步提升该类方法在理论研究和实际应用中的价值。
